% Archived from main.tex on 2026-10-16.
\section{Exposure to Extra-EU Import-Price Shocks}
\label{sec:app-import-shock-exposure}

\subsection{Direct and Indirect Effects in the Euro Area}

We apply the observed extra-EU import-price changes to imported final goods and intermediate inputs.
The FIGARO production network propagates the intermediate-input cost impulses across EA20 countries
and industries. We then map producer-price responses into consumption products and aggregate them
with the household expenditure weights used elsewhere in the paper. Throughout this section,
$k\in\{\text{energy},\text{food}\}$ indexes the import-price shock, $d$ denotes the destination
country, $i$ an imported product, $j$ a destination industry, and $g$ a household group. A
superscript $k$ labels the shock scenario; it is not an exponent.

Let $h_k$ denote the observed percentage change in the annual-average extra-EU import unit-value
index for shock $k$. For a final product $i$ absorbed by households in country $d$, let
$F_{di}^{\mathrm{extra}}$ denote imports from outside the EU27 and
$F_{di}^{\mathrm{all}}$ total household absorption, including domestic production and all imports.
The direct final-demand price impulse is
\begin{equation}
d_{di}^{k}=h_k\frac{F_{di}^{\mathrm{extra}}}{F_{di}^{\mathrm{all}}}.
\label{eq:direct-import-price-impulse}
\end{equation}
The ratio therefore scales the common external price shock by the extra-EU share of the product's
household supply. The vector $\mathbf d^k$ stacks $d_{di}^{k}$ over all country--product pairs.

For intermediate inputs, $\mathbf A$ is the FIGARO matrix of input coefficients for the EA20
country--industry network: each element records the amount of an input required per unit of gross
output in the purchasing industry. The vector $\mathbf b^k$ contains the first-round production-cost
impulses generated by extra-EU inputs affected by shock $k$. Its element for industry $j$ in country
$d$ is
\begin{equation}
b_{dj}^{k}
:=h_k\frac{\sum_{o\notin EU27}\sum_{i\in\mathcal S_k}Z_{oi,dj}}{x_{dj}},
\label{eq:main-import-cost-impulse}
\end{equation}
where $o$ indexes the origin country, $\mathcal S_k$ is the set of products affected by shock $k$,
$Z_{oi,dj}$ is the FIGARO intermediate flow of product $i$ from origin $o$ to industry $j$ in
destination country $d$, and $x_{dj}$ is that industry's gross output. Thus, $b_{dj}^{k}$ is the
first-round increase in unit production costs under complete pass-through of the import-price shock.
The vector $\mathbf b^k$ stacks these impulses over all EA20 country--industry pairs.

The corresponding producer-price response is
\begin{equation}
\Delta\mathbf p^{k}=(\mathbf I-\mathbf A^{\mathsf T})^{-1}\mathbf b^{k},
\label{eq:leontief-import-price-response}
\end{equation}
where $\mathbf I$ is the identity matrix and $(\mathbf I-\mathbf A^{\mathsf T})^{-1}$ is the
Leontief price inverse. The vector $\Delta\mathbf p^k$ includes both the initial imported-input cost
impulse and all subsequent rounds of cost propagation through domestic and intra-EA20 input links.

Finally, $\mathbf C$ is the normalized NACE--COICOP concordance matrix: multiplying an
industry-level price vector by $\mathbf C$ yields the corresponding vector of COICOP
consumption-price changes. The vector $\mathbf w_g$ contains the normalized COICOP expenditure
shares of household group $g$, so its elements sum to one. The total inflation exposure of group
$g$ is
\begin{equation}
\Delta\pi_g^k
:=\mathbf w_g^{\mathsf T}\mathbf C\mathbf d^{k}
+\mathbf w_g^{\mathsf T}\mathbf C\Delta\mathbf p^{k}.
\label{eq:import-shock-household-effect}
\end{equation}
The first term is the contribution of directly imported final goods; the second is the contribution
of imported intermediate inputs after production-network propagation. Because $h_k$ and all price
responses are expressed as percentage changes, $\Delta\pi_g^k$ is an inflation exposure in percentage
points. The effect on inflation inequality is
$\Delta\pi_{Q1-Q5}^k=\Delta\pi_{Q1}^k-\Delta\pi_{Q5}^k$, where a positive value means that shock $k$ raises
the inflation exposure of the first income quintile more than that of the fifth.

Panel (a) uses the observed increase in the annual-average extra-EU energy import unit-value index
from 2021 to 2022, and panel (b) uses the corresponding food import-price increase. The raw shocks
and their conversion into country--industry cost impulses are reported in
Appendix~\ref{sec:app-import-cost-impulses}. The shocks are common external-price innovations, while
bilateral import shares, production structures, and household baskets generate heterogeneous
country and household exposure.

%%%%%%%%%%%%%%%%%%%%%%%%%%%%%%%%%%%%%%%%%%%%%%%%%%%%%%%%%%%%%%
\begingroup
\renewcommand{\scriptsize}{\fontsize{8.5}{10}\selectfont}
\begin{table}[H]

\caption{\label{tab:shock-oil-gas-chain-ea20}Product-group contributions of observed extra-EU import-price shocks, euro area, 2022}
\centering
\small
\setlength{\tabcolsep}{3.0pt}
\begin{tabularx}{\textwidth}{Xrrrrrr}
\toprule
 & \multicolumn{3}{c}{Average inflation} & \multicolumn{3}{c}{Q1--Q5 gap}\\
\cmidrule(lr){2-4}\cmidrule(lr){5-7}
Product group & Total & Direct & Indirect & Total & Direct & Indirect\\
\midrule
\multicolumn{7}{l}{\textit{(a) Energy import-price shock (83.9\%)}}\\
Food and non-alcoholic beverages & 0.5 & 0.0 & 0.5 & 0.2 & 0.0 & 0.2\\
Alcoholic beverage, tobacco and narcotics & 0.1 & 0.0 & 0.1 & 0.1 & 0.0 & 0.1\\
Clothing and footwear & 0.1 & 0.0 & 0.1 & 0.0 & 0.0 & 0.0\\
Housing maintenance and furnishings & 0.2 & 0.0 & 0.2 & -0.1 & 0.0 & -0.1\\
Water, electricity, gas and other fuels & 1.3 & 0.4 & 0.9 & 1.6 & 0.5 & 1.1\\
Operation of personal transport equipment & 1.3 & 0.2 & 1.1 & 0.0 & 0.0 & 0.0\\
Transport (excluding vehicle operation) & 1.6 & 0.3 & 1.4 & -1.1 & -0.2 & -0.9\\
Other & 0.8 & 0.2 & 0.6 & -0.2 & -0.1 & -0.1\\
\textbf{Total} & \textbf{6.0} & \textbf{1.1} & \textbf{4.9} & \textbf{0.5} & \textbf{0.2} & \textbf{0.3}\\
\midrule
\multicolumn{7}{l}{\textit{(b) Food import-price shock (21.0\%)}}\\
Food and non-alcoholic beverages & 0.5 & 0.3 & 0.2 & 0.2 & 0.1 & 0.1\\
Alcoholic beverage, tobacco and narcotics & 0.1 & 0.1 & 0.0 & 0.1 & 0.0 & 0.0\\
Clothing and footwear & 0.0 & 0.0 & 0.0 & 0.0 & 0.0 & 0.0\\
Housing maintenance and furnishings & 0.0 & 0.0 & 0.0 & 0.0 & 0.0 & 0.0\\
Water, electricity, gas and other fuels & 0.0 & 0.0 & 0.0 & 0.0 & 0.0 & 0.0\\
Operation of personal transport equipment & 0.0 & 0.0 & 0.0 & 0.0 & 0.0 & 0.0\\
Transport (excluding vehicle operation) & 0.0 & 0.0 & 0.0 & 0.0 & 0.0 & 0.0\\
Other & 0.0 & 0.0 & 0.0 & 0.0 & 0.0 & 0.0\\
\textbf{Total} & \textbf{0.7} & \textbf{0.4} & \textbf{0.2} & \textbf{0.2} & \textbf{0.2} & \textbf{0.1}\\
\bottomrule
\end{tabularx}
\par\vspace{0.35em}\footnotesize\noindent Notes. Shocks are changes in the 2022 annual-average extra-EU27 import unit-value indices relative to 2021. Energy is SITC 3; food is SITC 0--1. Direct effects use extra-EU imports purchased by households. Indirect effects use extra-EU intermediate inputs and the full EA20 FIGARO propagation. Direct and indirect effects sum to the total before rounding. Sources: Eurostat international trade statistics, HBS, FIGARO, authors' calculations.
\end{table}

\endgroup
%%%%%%%%%%%%%%%%%%%%%%%%%%%%%%%%%%%%%%%%%%%%%%%%%%%%%%%%%%%%%%

Table~\ref{tab:shock-oil-gas-chain-ea20} shows that the observed energy import-price shock generates
6.0 percentage points of average EA20 non-rent inflation exposure: 1.1 points through imported
final goods and 4.9 points through intermediate inputs and production-network propagation.
Household energy, vehicle operation, and other transport contribute 1.3, 1.3, and 1.6 points,
respectively. The net Q1--Q5 effect is 0.5 point, comprising 0.2 point directly and 0.3 point
indirectly. Household energy contributes 1.6 points to the gap, partly offset by other transport
($-1.1$ points), housing maintenance, and other products.
The food import-price shock generates 0.7 point of average exposure and a Q1--Q5 gap of
0.2 point. Components and totals are rounded separately.

These are accounting exposures under complete pass-through, not causal estimates or a model fit to
observed HICP inflation. Retail margins, taxes, subsidies, regulated prices, substitution, and
incomplete pass-through can all attenuate or redistribute the observed consumer-price response.

\FloatBarrier

\subsection{Country Exposure}
\label{sec:app-import-cost-impulses}

For shock family $k$, let $h_k$ denote the observed change in the annual-average extra-EU27 import
unit-value index. The energy shock is 83.9 percent and the food shock is 21.0 percent between 2021
and 2022. These are common external-price innovations: country heterogeneity enters through import
exposure rather than through different values of $h_k$.

Equation~\eqref{eq:main-import-cost-impulse} converts the raw shock into an imported-input cost
impulse for each destination country and industry. An identical raw price shock therefore produces
different cost impulses across countries and industries because their extra-EU imported-input shares
differ.

The output-weighted country impulse reported in
Table~\ref{tab:extra-eu-import-cost-impulses-country-2022} is
\begin{equation}
\bar b_d^k
:=h_k\frac{\sum_j\sum_{o\notin EU27}\sum_{i\in\mathcal S_k}Z_{oi,dj}}
{\sum_j x_{dj}}.
\end{equation}
The euro-area row applies this ratio to EA20 sums; it is not a simple average of the country rows.

\begin{table}[H]

\caption{\label{tab:extra-eu-import-cost-impulses-country-2022}Import exposure and imported-input cost impulses by euro-area country, 2022}
\centering
\begin{tabular}[t]{l*{4}{>{\raggedleft\arraybackslash}p{1.45cm}}}
\toprule
 & \multicolumn{2}{c}{Energy products} & \multicolumn{2}{c}{Food products}\\
 & \multicolumn{2}{c}{Raw price shock: 83.9\%} & \multicolumn{2}{c}{Raw price shock: 21.0\%}\\
\cmidrule(lr){2-3}\cmidrule(lr){4-5}
Country & $M/x$ & $b$ & $M/x$ & $b$\\
\midrule
Austria & 0.79 & 0.66 & 0.11 & 0.02\\
Belgium & 2.36 & 1.98 & 0.24 & 0.05\\
Croatia & 2.62 & 2.20 & 0.24 & 0.05\\
Cyprus & 0.62 & 0.52 & 0.14 & 0.03\\
Estonia & 0.87 & 0.73 & 0.29 & 0.06\\
Finland & 2.27 & 1.91 & 0.13 & 0.03\\
France & 1.52 & 1.27 & 0.12 & 0.02\\
Germany & 1.20 & 1.01 & 0.10 & 0.02\\
Greece & 6.51 & 5.46 & 0.22 & 0.05\\
Ireland & 0.78 & 0.65 & 0.30 & 0.06\\
Italy & 2.00 & 1.68 & 0.29 & 0.06\\
Latvia & 0.73 & 0.61 & 0.23 & 0.05\\
Lithuania & 6.07 & 5.09 & 0.23 & 0.05\\
Luxembourg & 0.07 & 0.06 & 0.05 & 0.01\\
Malta & 0.81 & 0.68 & 0.10 & 0.02\\
Netherlands & 2.60 & 2.18 & 0.62 & 0.13\\
Portugal & 2.47 & 2.07 & 0.42 & 0.09\\
Slovakia & 2.96 & 2.48 & 0.16 & 0.03\\
Slovenia & 0.96 & 0.81 & 0.16 & 0.03\\
Spain & 2.70 & 2.27 & 0.29 & 0.06\\
\textbf{Euro area} & \textbf{1.80} & \textbf{1.51} & \textbf{0.21} & \textbf{0.04}\\
\bottomrule
\end{tabular}
\par\vspace{0.35em}\footnotesize\noindent Notes. The raw shocks are the common changes in the 2022 annual-average extra-EU27 import unit-value indices relative to 2021. $M/x$ is the extra-EU imported-input share of gross output, and $b=h(M/x)$ is the corresponding cost impulse; both columns are expressed in percent. Energy covers B and C19; food covers A01 and C10--12. The euro-area entries are ratios of EA20 sums. Sources: Eurostat international trade statistics, FIGARO, authors' calculations.
\end{table}


Table~\ref{tab:shock-extraeu-country-annual-2021-2023} repeats the simulation for each annual
energy import-price change from 2021 to 2023 and reports the resulting effects by country.
To isolate changes in the external shocks, the simulation holds the 2022 FIGARO network and all
2022 expenditure and country weights fixed across the three columns.

\begin{table}[H]

\caption{\label{tab:shock-extraeu-country-annual-2021-2023}Simulated country effects of the observed extra-EU energy import-price shock, 2021--2023}
\centering
\begin{tabular}[t]{l*{6}{>{\raggedleft\arraybackslash}p{1.10cm}}}
\toprule
 & \multicolumn{3}{c}{Effect on average inflation} & \multicolumn{3}{c}{Effect on inflation inequality}\\
\cmidrule(lr){2-4}\cmidrule(lr){5-7}
Country & 2021 & 2022 & 2023 & 2021 & 2022 & 2023\\
\midrule
Germany & 4.0 & 5.6 & -1.8 & 0.8 & 1.1 & -0.4\\
France & 5.1 & 7.1 & -2.3 & 0.8 & 1.1 & -0.4\\
Italy & 5.8 & 8.1 & -2.6 & 0.1 & 0.1 & -0.0\\
Spain & 4.3 & 5.9 & -1.9 & 0.4 & 0.6 & -0.2\\
Netherlands & 5.6 & 7.9 & -2.5 & -0.7 & -0.9 & 0.3\\
Belgium & 3.8 & 5.2 & -1.7 & 0.5 & 0.7 & -0.2\\
Austria & 2.6 & 3.5 & -1.1 & 0.1 & 0.2 & -0.1\\
Portugal & 4.9 & 6.8 & -2.2 & 0.7 & 1.0 & -0.3\\
Greece & 6.1 & 8.6 & -2.7 & 0.7 & 1.0 & -0.3\\
Finland & 3.1 & 4.4 & -1.4 & -0.9 & -1.3 & 0.4\\
Ireland & 3.8 & 5.3 & -1.7 & 2.4 & 3.3 & -1.0\\
Slovakia & 4.6 & 6.5 & -2.1 & 1.8 & 2.5 & -0.8\\
Croatia & 6.0 & 8.4 & -2.7 & 0.7 & 1.0 & -0.3\\
Lithuania & 6.2 & 8.6 & -2.8 & 0.0 & 0.0 & -0.0\\
Slovenia & 4.0 & 5.5 & -1.8 & -0.0 & -0.0 & 0.0\\
Luxembourg & 1.6 & 2.2 & -0.7 & 0.2 & 0.3 & -0.1\\
Latvia & 2.7 & 3.7 & -1.2 & 0.7 & 1.0 & -0.3\\
Estonia & 3.3 & 4.6 & -1.5 & 0.2 & 0.2 & -0.1\\
Cyprus & 4.1 & 5.7 & -1.8 & 1.6 & 2.2 & -0.7\\
Malta & 2.4 & 3.3 & -1.0 & 0.0 & 0.0 & -0.0\\
\textbf{Euro area} & \textbf{4.6} & \textbf{6.5} & \textbf{-2.1} & \textbf{0.5} & \textbf{0.7} & \textbf{-0.2}\\
\bottomrule
\end{tabular}
\par\vspace{0.35em}\footnotesize\noindent Notes. Entries are percentage-point effects of the energy import-price shock. Each column uses the change in the annual-average extra-EU27 energy import unit-value index relative to the previous year. The 2022 FIGARO network, product weights, household weights, and country weights are held fixed across years. Countries are ordered by their 2022 HICP country weight. Sources: Eurostat international trade statistics, HBS, FIGARO, authors' calculations.
\end{table}


\clearpage
